\chapter*{Preface}

I approach the preparation of this notebook with some ambivalence. The principal reason is that my doctoral research was conducted in control system, focusing primarily on robust state estimation and with additional work on adaptive control system. This background may give rise to at least the following two issues.

First, the distinction between a notebook that adheres strictly to its source materials and one that incorporates the author's (my) own interpretations may become blurred. Over the course of my research, I have developed personal understandings of, and in some cases broader philosophical perspectives on, many of the topics discussed in this notebook. Such views may inevitably invade the text even when this is not my conscious intention. A further concern is that many of these ideas have not been subjected to peer review, with the exception of material already published in my academic papers. Consequently, their correctness cannot be guaranteed. Indeed, some may prove inaccurate, particularly where they are overly speculative or ``visionary'' in nature.

Second, there is a risk that the notebook may become more extensive and detailed than its intended purpose requires. Ideally, it should remain clear, concise, and accessible. However, my tendency is to investigate a subject in increasing depth and to include more discussion and interpretation than may be strictly necessary. Although readers with a similar research background may find such material joyful to read, it may appear redundant or of limited practical value to a broader audience.

Nevertheless, notebooks in this series have already covered on a wide range of mathematical topics, and applied mathematics should also be represented. In my view, control systems constitute one of the most important applications of mathematics. It is therefore appropriate to include the subject here. The effectiveness of this approach will become apparent as the notebook develops.

There is one more thing worth mentioning. The readers may find overlaps among some of the chapters in this notebook. For example, parameter estimation and state estimation are introduced in control systems chapters, and they are introduced again with more details in their dedicated chapters in later part of the notebook. This is to make sure that each chapter is as much self-contained as possible. The readers shall feel free to skip repeated contents when reading this notebook.

\vspace{0.1in}
\noindent \rule{1in}{0.4pt}
\vspace{0.1in}

A list of key references of this notebook is given below. These materials are frequently cited in the notebook. For convenience, they are not given in the Reference section but listed here instead, whereas otherwise I will have to cite them everywhere in the text.

\vspace{0.1in}
\noindent \textbf{Control System}
\vspace{0.1in}

\begin{itemize}
	\item Åström, Karl Johan, and Richard Murray. Feedback systems: an introduction for scientists and engineers. Princeton university press, 2021.
	\item Åström, Karl J., and Björn Wittenmark. Computer-controlled systems: theory and design. Courier Corporation, 2013.
	\item Lewis, Frank L., Draguna Vrabie, and Vassilis L. Syrmos. Optimal control. John Wiley \& Sons, 2012.
	\item Rawlings, James B., David Q. Mayne, and Moritz M. Diehl. Model predictive control: theory, computation, and design, 2020.
	\item Åström, Karl J., and Björn Wittenmark. Adaptive Control. 2nd ed. Mineola, NY: Dover Publications, 2013.
\end{itemize}

\vspace{0.1in}
\noindent \textbf{System Identification}
\vspace{0.1in}

\begin{itemize}
	\item Hampel, Frank R., Elvezio M. Ronchetti, Peter J. Rousseeuw, and Werner A. Stahel. Robust Statistics: The Approach Based on Influence Functions. New York: Wiley, 1986.
	\item Lewis, Frank L., Lihua Xie, and Dan Popa. Optimal and robust estimation: with an introduction to stochastic control theory. CRC press, 2017.
\end{itemize}